\documentclass[11pt]{article}
\usepackage[utf8]{inputenc}
\usepackage[margin=1in]{geometry}
\usepackage{amsmath,amssymb}
\usepackage{booktabs}
\usepackage[colorlinks=true,linkcolor=blue,citecolor=blue,urlcolor=blue]{hyperref}
\usepackage{authblk}

% ============================================================================
% PAPER 3 (SHORT / WORKSHOP) -- practitioner reachability + compression guide.
% HONEST NOTE TO AUTHORS: this overlaps Papers 1-2. Submit as a short/workshop
% paper, OR fold into Paper 1 as a section, OR upgrade it into a SYSTEMS paper
% once real on-device (Jetson INT4 kernel) latency/energy numbers exist -- that
% hardware section is the content that would make this a standalone full paper.
% Placeholders below (Sec. 5) mark where on-device measurements go.
% ============================================================================

\title{\textbf{How Low Can You Go? A Reachability and Compression Guide\\
for Quantizing Vision--Language--Action Policies}}
\author[1]{Justin Williams}
\author[1]{Kishor Datta Gupta}
\author[1]{Roy George}
\author[1]{Mrinmoy Sarkar}
\affil[1]{Clark Atlanta University \quad \texttt{justinwilliamstech693@gmail.com}}
\date{}

\begin{document}
\maketitle

\begin{abstract}
Practitioners deploying Vision--Language--Action (VLA) policies on embedded robots need a
straight answer: \emph{how aggressively can I quantize before the policy stops working, and
how much smaller does that make the model?} We give an empirical, closed-loop answer for an
OpenVLA-OFT policy on LIBERO. We chart the reachable bit-width frontier---W4A4 and W3A8 are
``free'' ($\geq\!98\%$ task retention), W3A4 is the accuracy-trading cliff
($\sim\!10$--$15\%$ loss), and all sub-2-bit weight settings are unreachable regardless of
preconditioning---and pair it with an exact model-size accounting: W4A4 is $2.7\times$
smaller whole-model ($4\times$ on the $84\%$ quantized backbone), W3A8 is $3.16\times$.
We frame these as three deployment operating points and provide a decision rule.
\end{abstract}

\section{Introduction}
Quantization choices for VLA deployment are usually made by folklore. We replace folklore
with a measured map: closed-loop success and exact compression for every operating point,
on a real policy and benchmark \cite{openvlaoft,libero}.

\section{Setup}
OpenVLA-OFT policy (bf16 success $90.25\%$); closed-loop LIBERO, $400$ rollouts per
configuration; simulated quantization of the Llama backbone linears; weights per-channel
symmetric, activations per-token \cite{quarot}.

\section{The Reachability Frontier}
Table~\ref{tab:reach} summarizes. Three regimes appear cleanly.

\begin{table}[h]\centering
\caption{Reachability of quantization operating points (full eval where available;
otherwise \textsc{spatial} screen, lossless cluster $\approx 77$--$80\%$).}
\label{tab:reach}
\begin{tabular}{llcc}
\toprule
Regime & Configs & Success & Verdict \\
\midrule
Free lunch & W4, W4A8, \textbf{W4A4}, W3A8 & $88.8$--$90.5\%$ & deploy \\
Cliff edge & W3A4 (rotated) & $\sim\!70\%$ screen & trade-off only \\
Unreachable & any W2 (incl.\ W2A8 rotated) & $0\%$ & avoid \\
\bottomrule
\end{tabular}
\end{table}

The 2-bit wall is notable: rotation plus near-lossless 8-bit activations (W2A8) still yields
$0\%$, so 2-bit failure is a \emph{weight} representational limit, not an activation or
conditioning problem.

\section{Compression Accounting}
The checkpoint is $7.69$B parameters ($15.4$\,GB bf16); the quantized backbone is $6.48$B
($84\%$); the vision tower, embeddings, lm\_head, and action head ($1.22$B, $2.43$\,GB) stay
bf16. Table~\ref{tab:size} gives exact sizes.

\begin{table}[h]\centering
\caption{Whole-model and backbone compression by weight bit-width.}
\label{tab:size}
\begin{tabular}{lccc}
\toprule
Setting & Backbone & Whole model & Size \\
\midrule
bf16 & $1\times$ & $1\times$ & 15.4\,GB \\
W8 & $2.0\times$ & $1.73\times$ & 8.9\,GB \\
\textbf{W4} (W4A4) & $4.0\times$ & $2.71\times$ & 5.7\,GB \\
W3 (W3A8) & $5.3\times$ & $3.16\times$ & 4.9\,GB \\
\bottomrule
\end{tabular}
\end{table}

Activation bit-width does not change file size (activations are transient) but reduces
runtime activation bandwidth $\sim\!4\times$ at A4 and enables INT4 compute. Whole-model
compression is capped by the bf16 tower/heads; quantizing those is the obvious lever to
approach the backbone's $4\times$.

\section{On-Device Realization \emph{(to be completed)}}
\emph{This section is the path from a short paper to a systems paper.} We will report, for the
W4A4 + DCT operating point on an NVIDIA Jetson Orin Nano: end-to-end policy latency
(ms/query), peak memory, energy per rollout, and the maintained closed-loop success rate
under a real INT4 kernel. \textbf{[Placeholder: insert measured numbers.]}

\section{A Decision Rule}
\textbf{(1)} Default to \textbf{W4A4 + DCT}---near-lossless, $2.7\times$ smaller, $4\times$
activation bandwidth, any-dimension. \textbf{(2)} Need more compression and can spend $8$-bit
activations? Use \textbf{W3A8} ($3.16\times$). \textbf{(3)} Do \emph{not} attempt W2.
\textbf{(4)} Always apply a global orthogonal rotation before activation quantization; without
it, A4 collapses to $0\%$.

\section{Conclusion}
A measured reachability map removes the guesswork from VLA quantization: W4A4 is the deploy
point, W3A8 the compression point, W2 the wall, and the size wins are exact.

\begin{thebibliography}{99}
\small
\bibitem{openvlaoft} M.~Kim et al. ``Fine-Tuning Vision-Language-Action Models (OpenVLA-OFT).'' 2025.
\bibitem{libero} B.~Liu et al. ``LIBERO: Benchmarking Knowledge Transfer for Lifelong Robot Learning.'' 2023.
\bibitem{quarot} S.~Ashkboos et al. ``QuaRot: Outlier-Free 4-Bit Inference in Rotated LLMs.'' 2024.
\end{thebibliography}

\end{document}
